%%%%%%%%%%%%%%%%%%%%%%%%%%%%%%%%%%%%%%%%%%%%%%%%%%%%%%%%%%%%%%%%%%%%%%%%%%%%%%
%  qb-unit4.tex -- Unit IV: Web Design; Mobile Marketing
%%%%%%%%%%%%%%%%%%%%%%%%%%%%%%%%%%%%%%%%%%%%%%%%%%%%%%%%%%%%%%%%%%%%%%%%%%%%%%
\chapter{Web Design and Mobile Marketing}

\textit{Syllabus:} Web Design --- web design and optimisation of websites;
publishing a basic website; user-centered design and website optimisation;
design principles and website copy; website metrics and developing insight.
Mobile Marketing --- the difference between mobile advertising and mobile
marketing; using mobile marketing for sales promotions; online applications.

\textit{In the examination:} Part~B Questions~7 and~8 are set from this unit.
In both papers the pair mixed a web-design question with a mobile-marketing
question, so the two halves must be prepared together. Note that the 2025 paper
set Q7a on Google Analytics and Hotjar --- placed here because the blueprint
assigns Q7 to Unit~IV, and because the question is about developing
\emph{website insight}, which is a Unit~IV descriptor phrase.

\parta

\begin{qlist}

\qq{25-1.7}{Why is it important to integrate SEO in website development?}{SEE,
Jun/Jul 2025, Q1.7 --- 2 M}
Because the largest SEO factors are development decisions, not marketing
decisions: site architecture, URL structure, crawlability, rendering method,
page speed, mobile responsiveness, semantic markup and structured data are all
fixed while the site is being built. Integrating SEO from the start means the
site is indexable, fast and rank-ready on the day it launches, whereas
retrofitting it later means expensive rebuilds, URL migrations and the traffic
loss that accompanies them.
\scheme{1 M for the point that SEO elements are structural/development
decisions, 1 M for the consequence --- ranking from launch, and avoidance of
costly rework or traffic loss.}

\qq{25-1.8}{What is the role of geo-targeting in mobile marketing?}{SEE,
Jun/Jul 2025, Q1.8 --- 2 M}
Geo-targeting delivers advertisements, offers and content to users on the basis
of their physical location, determined by GPS, IP address, geofences or
in-store beacons. Its role is to make mobile messages contextually relevant ---
reaching a user when they are near a store, serving location-specific offers,
languages, stock and store details, and triggering proximity notifications ---
which raises response rates, eliminates spend on users outside the service
area, and produces measurable footfall.
\scheme{1 M for what geo-targeting is, 1 M for its role --- relevance,
proximity offers, reduced wastage or measurable footfall.}

\qq{24-1.10}{What is the primary difference between mobile advertising and
mobile marketing?}{SEE, Sept/Oct 2024, Q1.10; CIE-III, Jun 2026; Improvement
CIE, 4 Jun 2025; Improvement CIE, 28 Aug 2024 --- 2 M}
\textbf{Mobile advertising} is a paid, largely one-way promotional activity ---
in-app and in-game banners, interstitials, native and video ads, paid SMS and
paid search on mobile --- aimed at immediate reach, clicks and installs.
\textbf{Mobile marketing} is the whole strategy of reaching and serving
customers on mobile devices: responsive websites and mobile SEO, apps and push
notifications, SMS and WhatsApp, QR codes, location-based marketing, mobile
payments and mobile advertising itself. The primary difference is one of scope
--- advertising is a paid tactic contained within mobile marketing, which is
the broader relationship-building strategy.
\scheme{1 M for each definition, with the scope relationship --- advertising is
a subset of marketing --- carrying the distinction. A candidate who states only
``one is paid'' earns half.}

\qq{c4a1}{What is website optimization?}{CIE-III, Jun 2026 --- 2 M}
Website optimisation is the continuous, data-driven process of improving a
website's performance, speed, usability, content and search visibility so that
more of its visitors complete the actions the business wants. It covers
technical optimisation, content and copy, SEO and mobile experience, and it is
judged by conversion rate rather than by appearance.
\scheme{1 M for the process of improving performance, usability and
visibility, 1 M for the purpose --- more traffic and higher conversions.}

\qq{c4a2}{Mention any two principles of website design.}{CIE-III, Jun 2026 ---
2 M}
\textbf{Simplicity} --- remove everything that does not help the visitor
complete their task, since every additional element competes for the same
attention. \textbf{Consistency} --- the same layout, colours, typography,
button styles and terminology on every page, so the user learns the interface
once. Also acceptable: visual hierarchy, mobile responsiveness, accessibility,
fast loading, clear navigation.
\scheme{1 M each, with a phrase of justification.}

\qq{c4a3}{Mention two key principles of user-centered design in website
development.}{Improvement CIE, 4 Jun 2025 --- 2 M}
\textbf{Design from an explicit understanding of users, their tasks and their
environment} --- established through research rather than assumed.
\textbf{Iterative design refined through user involvement and user-centred
evaluation} --- prototypes tested with real users and improved in cycles, rather
than designed once and launched. Also acceptable: involving users throughout;
addressing the whole user experience; multidisciplinary design teams.
\scheme{1 M each.}

\qq{c4a4}{How can website metrics help in optimizing a website for marketing
effectiveness?}{Improvement CIE, 4 Jun 2025 --- 2 M}
Metrics locate problems precisely --- naming the page, step or device where
visitors are lost --- and rank potential fixes by likely impact, so effort goes
where the loss is largest. They also prove or disprove each change through
before-and-after measurement and A/B testing, reallocate budget towards the
channels that convert, and expose segment-specific failures that site-wide
averages conceal.
\scheme{1 M for diagnosis and prioritisation, 1 M for validation of changes or
budget reallocation.}

\qq{c4a5}{State two key elements that visitors commonly expect to see on a
well-optimized landing page.}{Improvement CIE, 4 Jun 2025 --- 2 M}
A \textbf{clear benefit-led headline} that matches the advertisement or link
which brought them there, so the visitor knows within seconds that they are in
the right place; and a \textbf{single prominent call-to-action button},
visually dominant and above the fold. Also acceptable: trust signals such as
ratings, testimonials or security badges; transparent pricing.
\scheme{1 M each.}

\qq{c4a6}{How does image optimization contribute to overall website
performance?}{Improvement CIE, 28 Aug 2024 --- 2 M}
Compressing images, serving modern formats and lazy-loading below-the-fold
assets reduces page weight and load time, which improves Core Web Vitals and
lowers bounce rate --- images are usually the largest component of page weight,
so this is the single highest-leverage speed fix. It also improves the mobile
experience on slow connections, and descriptive file names and alt text improve
both SEO and accessibility.
\scheme{1 M for speed, page weight and bounce rate, 1 M for the SEO,
accessibility or mobile benefit.}

\qq{c4a7}{List any four key metrics used to measure the success of a
website.}{Improvement CIE, 28 Aug 2024 --- 2 M}
\begin{apts}
\item Traffic --- users and sessions, split new against returning.
\item Bounce rate or engagement rate.
\item Average session duration or engagement time.
\item Conversion rate.
\end{apts}
Also acceptable: pages per session, page load time, traffic sources, exit rate,
cost per acquisition.
\scheme{Any four, half a mark each.}

\qq{c4a8}{How can push notifications be used in mobile marketing to enhance
sales promotions?}{Improvement CIE, 28 Aug 2024 --- 2 M}
Push notifications deliver time-critical offers instantly to the installed base
through an owned channel that no platform algorithm throttles. In sales
promotion they announce flash sales and create urgency with countdowns, recover
abandoned carts within minutes, trigger geo-fenced coupons when a user
approaches a store, and re-engage lapsed users with personalised offers --- each
deep-linked so the tap lands on the product rather than the home screen.
\scheme{1 M for the mechanism --- instant, owned, direct-to-device delivery ---
and 1 M for two or more promotional uses.}

\qq{c4a9}{How can user-friendly design impact the overall effectiveness of a
website?}{Improvement CIE, 28 Aug 2024 --- 2 M}
It reduces the cognitive effort needed to complete a task, so more visitors
finish what they came to do. The measurable effects are lower bounce rate,
longer engagement time, more pages per session and a higher conversion rate;
the indirect effects are greater trust in the brand, more return visits, and
better organic ranking, because search engines now reward the same engagement
behaviour that good usability produces.
\scheme{1 M for reduced effort and task completion, 1 M for the measurable
outcomes.}

\qq{c4a10}{List any two aspects of branding and consistency in web
design.}{SEE, Oct/Nov 2023 --- 2 M}
\textbf{Consistent visual identity} --- the same logo, colour palette,
typography, imagery style, iconography and button design applied on every page.
\textbf{Consistent tone of voice and terminology} --- in all copy, including
microcopy, error messages, form labels and CTA text, so the site reads as one
coherent personality rather than as several. Also acceptable: consistent
navigation and layout patterns across pages.
\scheme{1 M each.}

\qq{c4a11}{List any two differences between mobile web marketing and mobile app
marketing.}{SEE, Oct/Nov 2023 --- 2 M}
\textbf{Access and reach:} mobile web is reached through a browser with no
installation, so it has far broader reach and is better for acquisition; an app
requires download and installation, so it reaches a narrower but much more
engaged and higher-value audience. \textbf{Capability:} apps have full push
notification rights and device access --- camera, GPS, biometrics, wallet ---
and can work offline; mobile web has limited push capability and restricted
device access.
\scheme{1 M each.}

\qq{c4a12}{What is the impact of user behavior analysis on website
optimization?}{SEE, Oct/Nov 2023 --- 2 M}
It replaces assumption with observation: heatmaps, scroll maps, click maps and
session recordings show what users actually do --- which elements they look at,
which they ignore, where they hesitate, where they rage-click and how far down
the page they really scroll. Its impact is to generate the specific hypotheses
that A/B testing then validates, without which optimisation is guesswork
dressed as design opinion.
\scheme{1 M for what behaviour analysis reveals, 1 M for its role in generating
and validating optimisation hypotheses.}

\end{qlist}

\partb

\begin{qlist}

\qq{24-7a}{Outline the steps involved in publishing a basic website for a small
business.}{SEE, Sept/Oct 2024, Q7a --- 8 M}
\begin{apts}
\item \textbf{Define purpose, audience and required pages.} Decide what the
      site must achieve --- enquiries, orders, credibility, directions --- and
      list the pages that serve it: home, about, products or services,
      pricing, contact, and a blog if content is planned.
\item \textbf{Register a domain name.} Choose a short, memorable,
      brand-relevant name, prefer a \texttt{.com} or \texttt{.in} extension,
      and register it through a registrar with privacy protection and
      auto-renewal enabled.
\item \textbf{Choose hosting or a website builder.} Shared or cloud hosting
      with a CMS such as WordPress gives control and lower long-run cost; a
      hosted builder such as Wix, Squarespace or Shopify trades control for
      speed and simplicity. For a small business the decision turns on whether
      anyone will maintain it.
\item \textbf{Plan the structure and wireframe the pages.} Draw the navigation
      hierarchy and a simple wireframe of each page showing where the value
      proposition, images, trust signals and calls to action sit, before any
      design work begins.
\item \textbf{Select the platform, theme and design.} Choose a responsive,
      lightweight, mobile-first template; set brand colours, typography and
      logo; and keep the layout conventional --- a small business gains nothing
      from an experimental navigation pattern.
\item \textbf{Create the content.} Write clear customer-focused copy, take or
      source good photographs of the actual products or premises, and prepare
      product details, pricing, business hours, address, map and contact
      information.
\item \textbf{Build the pages and the conversion paths.} Assemble navigation,
      pages, contact and enquiry forms, WhatsApp or chat buttons, and prominent
      calls to action; verify responsiveness at mobile, tablet and desktop
      widths and check basic accessibility --- contrast, alt text, keyboard
      navigation.
\item \textbf{Add the essential integrations.} A working contact form with
      e-mail delivery, a payment gateway if selling, Google Maps, social
      profile links, and a claimed and completed Google Business Profile, which
      for a local business often delivers more traffic than the site itself.
\item \textbf{Apply on-page SEO basics.} Unique title tags and meta
      descriptions, one H1 per page, keyword-bearing readable URLs, image alt
      text, an XML sitemap and a robots.txt file.
\item \textbf{Secure the site.} Install an SSL certificate so the site serves
      over HTTPS, set up automated backups, keep the CMS and plugins updated,
      and add basic spam and firewall protection.
\item \textbf{Test before launch.} Check every page across browsers and
      devices, test all forms and payment flows end to end, run a page-speed
      test, and fix broken links and missing images.
\item \textbf{Publish and register with search engines.} Point the domain to
      the host, take the site live, then submit the sitemap through Google
      Search Console and install Google Analytics so measurement begins on day
      one.
\item \textbf{Maintain.} Monitor analytics and Search Console, refresh content,
      apply updates and backups on a schedule, and act on the enquiries the
      site produces --- a site nobody monitors is a cost, not an asset.
\end{apts}
\scheme{8 M. Roughly eight distinct steps at 1 M each, in a sensible order.
Domain, hosting, structure, content, build, SEO and security basics, testing,
go-live and measurement should all appear; sequence itself carries marks
because the question says ``outline the steps''.}

\qq{24-7b}{Discuss any four important website metrics.}{SEE, Sept/Oct 2024,
Q7b --- 8 M}
\begin{apts}
\item \textbf{Traffic --- users and sessions, new against returning.}
      \emph{What it is:} the number of distinct visitors and visits in a
      period. \emph{What it reveals:} the size and growth of the audience, and
      whether growth comes from acquisition (new users) or from loyalty
      (returning users). \emph{How to act on it:} a rising session count with a
      flat user count means the same people visiting more often --- an
      engagement success, not an acquisition one, and the two require different
      responses.
\item \textbf{Traffic sources and channels.} \emph{What it is:} the breakdown
      of arrivals into organic search, paid search, social, e-mail, referral
      and direct. \emph{What it reveals:} which investments actually produce
      visitors, and how dependent the business is on any single channel.
      \emph{How to act on it:} compare conversion rate and cost per acquisition
      by channel, not volume, and move budget accordingly; a channel producing
      40\% of traffic and 5\% of revenue is a problem, not an achievement.
\item \textbf{Bounce rate and engagement rate.} \emph{What it is:} the share of
      single-page, non-interacting sessions, and its inverse. \emph{What it
      reveals:} whether the landing experience matches the expectation that
      brought the visitor. \emph{How to act on it:} examine bounce rate by
      landing page and by source; a high figure on paid landing pages usually
      means the ad promised something the page does not deliver, and is fixed
      by message match, load speed and a clearer above-the-fold proposition.
\item \textbf{Average engagement time and pages per session.} \emph{What it
      is:} how long a visitor is actively engaged and how many pages they
      consume. \emph{What it reveals:} the depth of interest and the quality of
      internal linking and navigation. \emph{How to act on it:} low values on
      content pages point to weak or hard-to-read content; strengthening
      internal links and adding a clear next step raises both.
\item \textbf{Conversion rate and goal completions.} \emph{What it is:} the
      percentage of sessions that complete a defined valuable action.
      \emph{What it reveals:} whether the site does its job, and it is the only
      metric here that is directly a business outcome. \emph{How to act on it:}
      segment it by device, source and landing page to locate where conversion
      is lost, then A/B test the specific step.
\item \textbf{Page load time and Core Web Vitals.} \emph{What it is:} LCP, INP
      and CLS --- loading, responsiveness and visual stability. \emph{What it
      reveals:} technical experience quality, which affects both conversion and
      search ranking. \emph{How to act on it:} compress images, defer
      non-critical scripts, use a CDN and remove unused tags.
\item \textbf{Exit rate and top exit pages.} \emph{What it is:} the share of
      sessions ending on a given page. \emph{What it reveals:} where the
      journey breaks; a high exit rate on a checkout or form page is a defect,
      on a thank-you page it is expected.
\item \textbf{Cost per acquisition and return on ad spend.} \emph{What it is:}
      the marketing cost of each conversion and the revenue produced per unit
      of spend. \emph{What it reveals:} whether the whole operation is
      profitable, which no traffic metric can tell you.
\end{apts}
\scheme{8 M --- \emph{four} metrics at 2 M each, not eight at 1 M. Within each:
1 M for the definition and 1 M for its significance and how it is acted upon.
Answering with eight briefly named metrics is the commonest way to lose marks
on this question.}

\qq{24-8a}{Discuss the challenges of using mobile marketing for sales and
promotion.}{SEE, Sept/Oct 2024, Q8a --- 8 M}
\begin{apts}
\item \textbf{Screen size and attention constraints.} A promotion must
      communicate in a few square centimetres to a user who is walking,
      commuting or multitasking. Copy, imagery and the call to action all have
      to work at a glance, which makes complex offers very hard to convey.
\item \textbf{Device, operating-system and client fragmentation.} Thousands of
      device models, screen sizes, Android versions and e-mail or browser
      clients must all render the campaign correctly, multiplying design,
      development and testing cost.
\item \textbf{Privacy regulation and the loss of tracking.} GDPR, CCPA and the
      DPDP Act require consent for collection and use of personal and location
      data, while Apple's App Tracking Transparency and the deprecation of
      third-party identifiers have materially reduced targeting precision and
      attribution accuracy.
\item \textbf{Intrusiveness and permission fatigue.} The phone is a personal
      device, so poorly timed or excessive push notifications and SMS are
      experienced as an intrusion. The penalty is severe and immediate ---
      notification permissions revoked, the app uninstalled, the number
      blocked.
\item \textbf{Ad avoidance and quality problems.} Ad blockers, banner
      blindness, accidental ``fat-finger'' clicks that inflate reported
      performance, and mobile ad fraud through click farms and install fraud
      all degrade the value of the impressions bought.
\item \textbf{High acquisition cost and poor app retention.} Cost per install
      continues to rise, while a large majority of installed apps are abandoned
      within the first week --- so the promotion pays for an install that
      produces no second visit.
\item \textbf{Network, data cost and device constraints.} Heavy creatives,
      video and large app updates fail on weak connections or are declined by
      users conscious of data cost and battery, particularly in price-sensitive
      markets.
\item \textbf{Measurement and attribution complexity.} Journeys cross app, web
      and physical store on the same device and across devices; without
      deterministic identity resolution, credit for a sale cannot be assigned
      reliably, which makes budget defence difficult.
\item \textbf{Security and trust.} SMS phishing, malicious apps and payment
      fraud have made users cautious about links and permissions, so legitimate
      promotional messages are treated with suspicion.
\item \textbf{Platform dependence and rapid change.} App store policies,
      operating-system releases and platform algorithm changes can invalidate a
      working tactic without notice, forcing continuous re-investment.
\item \textbf{Content, language and localisation demands.} Mobile audiences are
      linguistically diverse, and campaigns often require localisation into
      several languages and cultural contexts to work at all.
\item \textbf{Clutter and competition.} Every brand is targeting the same
      screen, so the marginal value of one more notification falls while the
      cost of standing out rises.
\end{apts}
\scheme{8 M for eight challenges at 1 M each, named and explained. Marks are
for distinct challenges --- ``small screen'' and ``limited attention'' will be
read as one point, not two.}

\qq{24-8b}{Elucidate the importance of understanding user-centered design
concepts while developing digital marketing strategies.}{SEE, Sept/Oct 2024,
Q8b --- 8 M}
\textbf{User-centered design (UCD)} is an iterative design process in which the
needs, goals, contexts and limitations of the actual user drive every decision
at every stage. Its cycle is: research the users \ra\ specify the context and
requirements \ra\ design a solution \ra\ evaluate it against real users \ra\
repeat. Its instruments are personas, journey maps, wireframes, prototypes,
usability testing and accessibility standards.

\begin{apts}
\item \textbf{It grounds strategy in evidence rather than assumption.} UCD
      research replaces what the marketing team believes customers want with
      what observed users actually do, which is the single most common source
      of wasted digital spend.
\item \textbf{It removes friction and so raises conversion.} Marketing pays to
      bring a visitor to a destination; if that destination is confusing, slow
      or demands too much, the acquisition spend is lost at the last step. UCD
      is what makes the traffic convert.
\item \textbf{It produces the personas and journey maps the strategy runs on.}
      The same research output --- who the user is, what they are trying to do,
      where they struggle --- drives targeting, channel selection, messaging
      and content sequencing, so UCD and marketing planning share a foundation.
\item \textbf{It builds trust and credibility.} Clarity, honest labelling,
      predictable navigation, visible pricing and transparent policies are
      design decisions that determine whether a first-time visitor believes the
      brand --- and trust is a precondition of conversion, not a consequence.
\item \textbf{It makes the offering accessible and inclusive.} Designing to
      WCAG standards --- contrast, text sizing, keyboard operability,
      screen-reader compatibility, captions --- extends reach to users with
      disabilities and to everyone in poor conditions, and reduces legal risk.
\item \textbf{It enforces context realism, especially mobile.} UCD asks where
      and how the product is actually used --- one-handed, on mobile data, in
      sunlight, between meetings --- and produces mobile-first, low-bandwidth,
      short-attention designs that match reality instead of the studio.
\item \textbf{It reduces cost and rework.} A usability problem found in a
      paper prototype costs almost nothing; the same problem found after
      launch costs a redevelopment cycle and the revenue lost meanwhile.
      Iterative evaluation moves discovery earlier, where it is cheap.
\item \textbf{It supports SEO and engagement indirectly.} Clear structure, fast
      loading, readable content and satisfied users produce the engagement
      signals and dwell times that search engines reward, so good UX and good
      organic visibility reinforce each other.
\item \textbf{It creates loyalty and advocacy.} A pleasant, low-effort
      experience produces repeat visits, positive reviews and referrals, which
      lower acquisition cost and raise lifetime value --- the compounding
      returns that campaign spending alone never generates.
\item \textbf{It gives the strategy a testable, adaptive method.} Usability
      testing and A/B testing keep the strategy responsive to changing user
      behaviour, so decisions can be revisited on evidence rather than defended
      on authority.
\end{apts}

\textbf{Conclusion.} Digital marketing controls the arrival; user-centered
design controls what happens after it. Since the return on every rupee spent on
acquisition is realised only at the destination, understanding UCD is not a
design specialism adjacent to marketing strategy --- it is the condition on
which that strategy's numbers depend.
\scheme{8 M. 2 M for defining UCD and its process, 5 M for five or six
substantiated reasons at roughly 1 M each, 1 M for a conclusion connecting UCD
to marketing outcomes.}

\qq{25-7a}{Demonstrate how tools like Google Analytics and Hotjar can be used
to develop insights into user behavior and improve marketing decisions.}{SEE,
Jun/Jul 2025, Q7a --- 8 M}
The two tools answer different questions and are used together. \textbf{Google
Analytics is quantitative} --- it tells you \emph{what} happened, to how many
people, and how much it cost. \textbf{Hotjar is qualitative} --- it shows
\emph{why} it happened.

\textit{Google Analytics --- the what}
\begin{apts}
\item \textbf{Acquisition analysis.} Sessions, engaged sessions, conversions
      and revenue by channel, source and campaign, showing which marketing
      investments produce buyers rather than visitors.
\item \textbf{Landing-page and content performance.} Entrances, engagement
      rate, average engagement time and exits per page, identifying which
      content earns attention and which loses it.
\item \textbf{Funnel and path exploration.} Step-by-step drop-off through the
      conversion funnel, and the actual paths users take, which are frequently
      not the paths the site was designed around.
\item \textbf{Segmentation.} The same funnel split by device, geography, new
      against returning and traffic source --- which is usually where the real
      finding is, for example a mobile conversion rate one-third of desktop.
\item \textbf{Key events and value.} Conversion tracking with assigned values,
      cost per acquisition and ROAS, so behaviour is translated into money.
\end{apts}

\textit{Hotjar --- the why}
\begin{apts}\setcounter{enumi}{5}
\item \textbf{Heatmaps.} Click, move and scroll maps show what users actually
      look at and interact with --- typically revealing that the key call to
      action sits below the point 60\% of users ever scroll to, or that users
      repeatedly click a non-clickable image.
\item \textbf{Session recordings.} Replays of individual sessions expose rage
      clicks, hesitation, repeated back-navigation and the exact field at which
      a form is abandoned.
\item \textbf{Form and funnel analytics.} Field-level analysis naming the input
      that causes abandonment --- typically a phone number, a mandatory GST or
      PAN field, or an unexplained password rule.
\item \textbf{On-site polls, surveys and feedback widgets.} Voice-of-customer
      data captured in the moment --- ``what stopped you from completing your
      purchase today?'' --- which no behavioural dataset can supply.
\end{apts}

\textit{The combined workflow, demonstrated}
\begin{apts}\setcounter{enumi}{9}
\item GA4 reports that mobile checkout completion is 22\% against 61\% on
      desktop --- a quantified problem, but no cause. Hotjar recordings of
      mobile checkout sessions show users reaching the coupon field, leaving to
      search for a code, and never returning; the heatmap confirms heavy
      tapping on that field. A hypothesis follows: the coupon field is causing
      abandonment. The remedy --- collapse it behind a small ``have a code?''
      link --- is A/B tested, and GA4 measures the resulting lift. That loop,
      \emph{quantify \ra\ explain \ra\ hypothesise \ra\ test \ra\ re-measure},
      is the demonstration the question asks for.
\end{apts}

\textit{Marketing decisions improved:} reallocation of budget to the channels
GA4 shows convert; redesign of the landing pages Hotjar shows confuse;
shortening of forms; reordering of page content to match real scroll depth;
mobile-first prioritisation; and content commissioned against the questions
surveys and site-search reveal.
\scheme{8 M. 3 M for Google Analytics capabilities, 3 M for Hotjar
capabilities, 2 M for the combined workflow and the marketing decisions it
changes. The command word is ``demonstrate'', so a worked example carries
marks --- listing features of both tools without showing them used together
answers only part of the question.}

\qq{25-7b}{Formulate a mobile-based marketing plan for a local food delivery
startup to reach a targeted student segment.}{SEE, Jun/Jul 2025, Q7b --- 8 M}
\begin{apts}
\item \textbf{Objective.} 5{,}000 app installs and 1{,}500 monthly orders from
      three named campuses within six months, at a customer acquisition cost
      below \Rs 120 and a repeat-order rate above 40\%.
\item \textbf{Segment insight.} Students aged 18--24 living in hostels and
      shared flats: highly price-sensitive, order late at night and during exam
      weeks, order in groups to split delivery fees, pay by UPI, discover
      through Instagram and WhatsApp rather than search, and are strongly
      influenced by peers and by referral credits.
\item \textbf{Positioning and offer design.} Student-priced combos under
      \Rs 99; free delivery on verification of a college ID; a group-order
      feature that splits the bill; a midnight menu during examination periods;
      and a referral scheme giving credit to both parties, which is the
      cheapest acquisition channel in a dense campus network.
\item \textbf{App and mobile experience.} A lightweight app that installs
      quickly on mid-range phones and poor connections; OTP sign-up in one
      step; saved hostel and block addresses; one-tap reorder; live tracking;
      UPI as the default payment; and a fast progressive web app for students
      unwilling to install anything. App-store optimisation on ``student food
      delivery \emph{city}'' and hostel-related terms.
\item \textbf{Channel mix.} Geofenced Instagram and Snapchat advertising drawn
      tightly around campus and hostel blocks; click-to-WhatsApp ads for
      ordering without an install; a campus ambassador programme paying
      commission per referred order; QR-code posters in canteens, hostel notice
      boards and library entrances; and opt-in WhatsApp broadcast lists per
      campus.
\item \textbf{Location and timing tactics.} Geofencing and, where available,
      beacons to trigger proximity offers; dayparting concentrated on the
      lunch break, the 8--10 pm window and the midnight peak; weather-triggered
      promotions on rainy evenings; and calendar-triggered campaigns keyed to
      examination and fest schedules.
\item \textbf{Retention programme.} Loyalty points and order streaks;
      personalised push notifications based on order history, capped in
      frequency to protect the notification permission; a subscription pass
      offering free delivery for a month; and win-back offers to users inactive
      for 21 days.
\item \textbf{Measurement and optimisation.} Track installs, cost per install
      and CAC, notification opt-in and open rates, order conversion rate,
      average order value, repeat-order rate, Day-7 and Day-30 retention and
      ROAS in Firebase and GA4 with a mobile measurement partner for
      attribution; run weekly creative A/B tests; review campus-level unit
      economics monthly and stop what does not pay.
\end{apts}
\scheme{8 M for eight plan components at 1 M each. The plan must be visibly
\emph{mobile-based} --- app, push, WhatsApp, geofencing, QR --- and visibly
aimed at \emph{students}. A general digital marketing plan, however good, will
lose the application marks that this question is mostly testing.}

\qq{25-8a}{Evaluate the effectiveness of mobile marketing in facing market
changes due to rapid technological advancements.}{SEE, Jun/Jul 2025, Q8a ---
8 M}
\textit{The case for effectiveness}
\begin{apts}
\item \textbf{Ubiquity and permanence of the channel.} The smartphone is now
      the primary internet device for most users and is carried at all times,
      so mobile marketing reaches the customer wherever the market moves --- it
      does not depend on the customer coming to a desk, a shop or a television.
\item \textbf{Speed of response.} Push notifications, in-app messages, SMS and
      WhatsApp campaigns can be conceived and delivered within hours, so a
      brand can react to a competitor's move, a price change, a stock-out or a
      sudden demand shift almost immediately --- something print, television
      and even e-mail cannot match.
\item \textbf{Data richness and personalisation.} Device, location, app usage
      and purchase history support individually tailored offers and
      AI-driven recommendation, so the offering can be re-tuned continuously as
      preferences shift.
\item \textbf{Context awareness.} Geofencing, proximity beacons and real-time
      triggers allow marketing that responds to where and when the customer
      actually is --- a capability no traditional medium possesses at all.
\item \textbf{Measurement and iteration speed.} Attribution, cohort retention
      and in-app analytics close the feedback loop in days rather than
      quarters, so strategy adapts at the speed the technology changes.
\item \textbf{Native absorption of new technology.} Each wave of advancement
      has arrived on the phone first and been absorbed by mobile marketing
      quickly --- 5G making rich video viable, AR try-on, voice search and
      assistants, conversational commerce through chatbots, UPI and one-tap
      payment, wearables and connected devices. The channel has repeatedly
      demonstrated it can adopt rather than be displaced.
\item \textbf{Cost-efficiency and accessibility.} Precise targeting and
      low production cost make it viable for small firms, which means the
      channel spreads adaptation across the whole market rather than only to
      large incumbents.
\end{apts}

\textit{The limits on that effectiveness}
\begin{apts}\setcounter{enumi}{7}
\item \textbf{Privacy regulation and identifier loss} --- ATT, cookie
      deprecation and consent requirements have removed much of the targeting
      precision on which the personalisation argument rested.
\item \textbf{Saturation and intrusion} --- notification fatigue, ad blocking
      and permission revocation mean the same immediacy that makes mobile
      responsive also makes it easy to alienate users.
\item \textbf{Fragmentation, obsolescence and cost} --- constant device, OS and
      platform change forces continuous re-investment, and high app acquisition
      costs with poor retention undermine the cost argument.
\item \textbf{Unequal access} --- device capability, data cost and network
      quality vary widely, so a strategy assuming high-end phones excludes part
      of the market.
\end{apts}

\textbf{Evaluation.} On balance mobile marketing is the most adaptive channel
currently available: its speed, data and context-awareness let a business
absorb technological change rather than be overtaken by it, and its record of
assimilating each successive advancement supports that. But its effectiveness
is conditional, not automatic. It holds only where the brand builds
consent-based first-party data to replace lost third-party targeting, restrains
frequency so that immediacy does not become intrusion, and continues investing
in experience as devices change. Where those conditions fail, the channel's
greatest strength --- direct access to a personal device --- becomes its
greatest liability.
\scheme{8 M. 4 M for the effectiveness arguments, 2 M for the limitations or
challenges, 2 M for a reasoned overall judgement. ``Evaluate'' requires the
verdict; an answer that lists advantages only is incomplete.}

\qq{25-8b}{Create a strategy to use new digital channels like mobile apps,
chatbots, and social media integration to strengthen a brand's global online
presence.}{SEE, Jun/Jul 2025, Q8b --- 8 M}
\begin{apts}
\item \textbf{Strategic objective and the glocal principle.} One consistent
      global brand identity, promise and visual system, executed with local
      relevance in language, culture, payment and festival calendar. The
      governing rule is ``global framework, local expression'', and success is
      measured by global share of digital voice, engagement rate, assisted
      revenue and consistency of brand perception across markets.
\item \textbf{The mobile app as the owned hub.} Build the app as the one
      channel the brand controls entirely: multilingual and multi-currency,
      region-aware content and inventory, a loyalty programme that works across
      markets, offline-capable for low-connectivity regions, and permissioned
      push notifications scheduled to local time zones. Support it with
      app-store optimisation in each market's language and a progressive web
      app for regions where installs are resisted.
\item \textbf{Chatbots for round-the-clock multilingual service.} Deploy
      conversational agents on the website, in the app, and on WhatsApp,
      Messenger and Instagram DMs, handling FAQs, order tracking, product
      finding, appointment booking and lead qualification in the local
      language, with clean escalation to a human agent. For a global business
      this is what makes 24-hour presence economically possible across time
      zones, and it captures structured intent data as a by-product.
\item \textbf{Social media integration.} Operate regional handles under a
      shared global content template and calendar; enable social login and
      social commerce with shoppable posts and in-platform checkout; run
      user-generated-content and hashtag campaigns that travel across markets;
      and build a per-market influencer network. Platform selection follows the
      market, not headquarters --- WeChat and Douyin in China, LINE in Japan,
      WhatsApp in India and Brazil, TikTok, Instagram and LinkedIn elsewhere.
\item \textbf{Data and technology integration.} Join app, web, chatbot, social
      and CRM identities in a customer data platform, with server-side tagging
      through GTM feeding one analytics stack. Without this the channels remain
      three separate silos; with it, a customer recognised on WhatsApp in the
      morning receives a consistent experience in the app that evening. This
      integration layer is what converts a list of channels into a strategy.
\item \textbf{Content and localisation.} Transcreate rather than translate ---
      adapt idiom, imagery, models, colour associations and humour to each
      market; localise pricing, payment methods and delivery promises; comply
      with local advertising and product-claim regulation; and meet
      accessibility standards so the presence is genuinely global rather than
      merely multinational.
\item \textbf{Emerging-channel layer.} Extend deliberately into AR try-on and
      virtual product experiences, voice assistants and voice search
      optimisation, short-form video, community platforms such as Discord and
      Telegram, and connected TV --- piloting each in one market before global
      rollout rather than launching everywhere at once.
\item \textbf{Governance, compliance and measurement.} Publish global brand and
      tone guidelines with a regional approval workflow; run a consent
      management platform satisfying GDPR, CCPA and DPDP simultaneously;
      maintain a crisis-response playbook covering all time zones; and report a
      single global dashboard --- reach, engagement rate, app monthly active
      users, chatbot containment rate and CSAT, social commerce revenue,
      conversion and ROAS --- reviewed monthly, with a test-and-learn pilot
      protocol before any market-wide launch.
\end{apts}
\scheme{8 M for eight strategy components at 1 M each. All three named channels
--- app, chatbot, social integration --- must be addressed substantively, and
the answer should show them \emph{integrated} rather than described separately;
the data and governance layers are what distinguish a strategy from a list.}

\qq{c4b1}{Explain the concept of web design and the importance of website
optimization.}{CIE-III, Jun 2026 --- 10 M}
\textbf{Web design} is the planning and creation of a website's structure,
appearance and interaction --- layout and grid, colour palette, typography,
imagery and graphics, navigation architecture, responsiveness across devices,
and the interaction patterns that let a user complete a task. It is the
discipline in which visual communication and usability meet business objective.

\textit{Types of optimisation}
\begin{apts}
\item \textbf{Technical} --- load speed, Core Web Vitals, clean code, uptime,
      security and HTTPS.
\item \textbf{Content} --- clarity, structure, readability, relevance and
      persuasive copy.
\item \textbf{SEO} --- crawlability, metadata, heading structure, internal
      linking, structured data.
\item \textbf{Mobile} --- responsive layout, tap targets, performance on weak
      connections.
\end{apts}

\textit{Importance of website optimisation}
\begin{apts}\setcounter{enumi}{4}
\item \textbf{It raises conversions from existing traffic,} which is almost
      always cheaper than buying more traffic.
\item \textbf{It improves user experience and satisfaction,} directly lowering
      bounce and abandonment.
\item \textbf{It improves search visibility,} because speed, mobile-friendliness
      and structure are ranking factors as well as usability factors.
\item \textbf{It reduces cost per acquisition across every channel at once,}
      since every channel points at the same destination.
\item \textbf{It builds credibility and trust} --- a slow, broken or confusing
      site is read as an unreliable business.
\item \textbf{It sustains competitiveness,} because user expectations, devices
      and algorithms move continuously, so an unoptimised site degrades even
      when nothing about it changes.
\end{apts}

\textbf{Benefits in summary:} more traffic, better engagement and more revenue
from the same marketing spend.
\scheme{10 M. Roughly 4 M for the concept of web design and its elements, 2 M
for the types of optimisation, and 4 M for the importance points. Both halves of
the question are named explicitly and must both appear.}

\qq{c4b2}{A start-up notices that visitors leave their website within a few
seconds. As a web designer, explain how you would optimize the website to
improve user engagement and conversion rate.}{CIE-III, Jun 2026 --- 10 M}
\begin{apts}
\item \textbf{Diagnose before redesigning.} Read GA4 for bounce by landing page,
      device and source, and watch Hotjar recordings and scroll maps to see
      \emph{why} people leave. Redesigning without diagnosis usually moves the
      problem rather than fixing it.
\item \textbf{Fix load speed first.} Visitors leaving within seconds are often
      leaving before the page renders. Compress and lazy-load images, serve
      modern formats, defer non-critical scripts, remove unused tags, enable
      caching and a CDN, and target a load under three seconds.
\item \textbf{Make it genuinely mobile responsive.} Most traffic is mobile;
      test real breakpoints, ensure tap targets of at least 44 pixels, type of
      16 pixels or more, and no horizontal scrolling.
\item \textbf{Rewrite the above-the-fold section.} One benefit-led headline
      stating what the business does for the visitor, a supporting sub-line, one
      dominant call to action --- the visitor must know within five seconds that
      they are in the right place.
\item \textbf{Enforce message match.} The landing page must repeat the promise
      of the advertisement, keyword or link that produced the visit; mismatch is
      the commonest cause of instant exits.
\item \textbf{Simplify navigation and structure.} A clear menu, logical
      hierarchy, working search, breadcrumbs, and no more than a handful of
      top-level choices.
\item \textbf{Improve content quality and scannability.} Short paragraphs,
      descriptive subheadings, bullets, supporting visuals, and the answer to
      the visitor's question near the top rather than after three screens of
      company history.
\item \textbf{Strengthen calls to action and reduce friction.} One visually
      dominant accent-coloured button with action-verb copy, repeated at
      intervals; forms cut to the minimum fields; guest checkout; visible
      progress indicators.
\item \textbf{Add trust signals.} Reviews and ratings, testimonials, security
      badges, clear contact details, refund and privacy policies --- a new
      start-up has no brand recognition, so trust must be constructed on the
      page.
\item \textbf{Validate and iterate.} A/B test the headline, hero image and CTA
      one change at a time; monitor bounce rate, engagement time, scroll depth
      and conversion rate weekly; and repeat the cycle rather than treating the
      fix as final.
\end{apts}
\scheme{10 M for ten optimisation actions at 1 M each. Opening with diagnosis
and closing with A/B validation are the two moves that lift this above a generic
checklist --- the question says ``as a web designer'', so a professional method
is being examined, not a list of tips.}

\qq{c4b3}{A college plans to publish its departmental website for admissions and
student activities. Explain the complete process involved in designing and
publishing the website.}{CIE-III, Jun 2026 --- 10 M}
\begin{apts}
\item \textbf{Requirement analysis.} Establish the objectives --- admissions
      enquiries, student information, activity updates, placement data --- and
      the audiences: prospective students and parents, current students, faculty
      and recruiters. Each wants something different from the same site.
\item \textbf{Information architecture and website design.} Plan the sitemap
      (home, about, programmes, admissions, faculty, research, activities,
      placements, contact), wireframe the key pages, and design the visual
      layout within the institution's brand identity.
\item \textbf{Content development.} Write and collect programme details,
      eligibility and fee structure, faculty profiles, syllabi, event
      calendars, photographs, accreditation and placement statistics --- for an
      admissions site, content accuracy is a compliance matter, not only a
      quality one.
\item \textbf{Domain registration.} Register an appropriate domain, ideally
      under the institution's existing \texttt{.edu.in} or \texttt{.ac.in}
      domain, which carries credibility that a generic domain does not.
\item \textbf{Web hosting.} Choose hosting sized for admission-season traffic
      spikes, with adequate bandwidth, uptime guarantees, automated backups and
      an SSL certificate.
\item \textbf{Website development.} Build the pages on a CMS such as WordPress
      or Drupal so that non-technical staff can update notices; implement
      responsive layouts, the enquiry and application forms, the notice board
      and the events calendar.
\item \textbf{Testing.} Verify across browsers and devices; test every form and
      the application workflow end to end; check load speed, broken links and
      accessibility --- a public institution's site has an accessibility
      obligation; and have faculty proof the academic content.
\item \textbf{Pre-launch technical setup.} Install SSL, generate and submit the
      XML sitemap, configure robots.txt, add title tags and meta descriptions,
      and install Google Analytics and Search Console so measurement starts on
      day one.
\item \textbf{Publishing.} Point the domain to the host, deploy, verify the live
      site, and announce it through the institution's existing channels ---
      notice boards, social handles, admission communications.
\item \textbf{Maintenance.} Assign an owner; update notices, results and event
      information on a schedule; refresh admission content each cycle; apply
      security patches and backups; and review analytics to see which
      information prospective students actually look for --- an unmaintained
      college site loses credibility faster than no site at all.
\end{apts}
\scheme{10 M for the ten stages at 1 M each, in correct sequence. The official
scheme names requirement analysis, design, content development, domain
registration, hosting, development, testing, publishing and maintenance ---
covering those nine secures the marks; the pre-launch technical setup and the
college-specific detail earn the top band.}

\qq{c4b4}{Discuss what website visitors typically look for when they land on a
business webpage, and suggest design and content strategies to meet those
expectations.}{Improvement CIE, 4 Jun 2025 --- 10 M}
Answer as visitor question paired with delivery strategy --- that pairing is
what the scheme rewards.
\begin{apts}
\item \textbf{``Am I in the right place?''} \emph{Strategy:} a headline stating
      plainly what the business does and for whom, matching the source that
      brought them, visible without scrolling.
\item \textbf{``What exactly do you offer?''} \emph{Strategy:} a clear product
      or service overview with benefit-led descriptions, images and categories
      one click from the homepage.
\item \textbf{``What does it cost?''} \emph{Strategy:} transparent pricing, or a
      clear explanation of how pricing works when it must be quoted. Hidden
      pricing is the most common cause of exit on a business site.
\item \textbf{``Can I trust you?''} \emph{Strategy:} reviews and ratings,
      testimonials with names and photographs, client logos, certifications,
      case studies, a real address and phone number, and an About page with
      actual people.
\item \textbf{``How do I contact you?''} \emph{Strategy:} contact details in the
      header and footer, a short enquiry form, WhatsApp or chat, a map, and
      stated response times.
\item \textbf{``What do I do next?''} \emph{Strategy:} one obvious, dominant
      call to action per page with action-verb copy, repeated at intervals down
      the page.
\item \textbf{``Is this current and maintained?''} \emph{Strategy:} recent
      updates, current dates, working links, live stock or availability --- an
      obviously stale site reads as a closed business.
\item \textbf{``Will this work on my phone, and quickly?''}
      \emph{Strategy:} mobile-first responsive design, sub-three-second load,
      compressed images, readable type and thumb-sized targets.
\item \textbf{``What are the terms if something goes wrong?''}
      \emph{Strategy:} visible shipping, returns, warranty, cancellation and
      privacy policies --- risk reversal converts hesitant buyers.
\item \textbf{Underlying design and copy principles.} Simplicity and generous
      whitespace; clear visual hierarchy putting the most important thing
      highest; consistency of colour, type and button style; accessibility of
      contrast and alt text; and copy that is customer-focused, specific,
      scannable and free of jargon --- ``you'' language rather than ``we''
      language.
\end{apts}
\scheme{10 M. Roughly 6--7 M for the visitor expectations each paired with a
strategy, and 3--4 M for the general design and copy principles. Listing what
visitors want without saying how the site delivers it answers only half the
question.}

\qq{c4b5}{Design a user-centered website layout for a new e-commerce startup
targeting young urban professionals. How would you ensure optimal user
experience and engagement? Include key design principles, optimization
strategies and website copy suggestions.}{Improvement CIE, 4 Jun 2025 --- 10 M}
\textbf{Research first.} Persona: 25--35, urban, mobile-first, time-poor,
price-aware but convenience-driven. Primary task: find and buy quickly on a
phone, often during a commute. Every decision below follows from that.

\begin{apts}
\item \textbf{Layout --- above the fold.} Sticky header carrying logo, prominent
      search bar, account and cart; a hero section with one benefit headline,
      one supporting line and one call to action.
\item \textbf{Layout --- body.} A three-column category grid (single column on
      mobile); a social-proof band with aggregate ratings and review count; a
      personalised ``recommended for you'' row; a delivery-promise strip; and a
      simple footer carrying trust, policy and contact links.
\item \textbf{Design principle --- mobile-first responsive.} Designed at the
      phone breakpoint first and expanded upward, not the reverse.
\item \textbf{Design principle --- visual hierarchy.} Search bar and primary CTA
      given the greatest visual weight, because search is how a time-poor buyer
      shops.
\item \textbf{Design principle --- simplicity and consistency.} Generous
      whitespace, one accent colour reserved exclusively for actions, one
      button style, type at 16 pixels or larger, tap targets of at least 44
      pixels, and accessible contrast with alt text throughout.
\item \textbf{Optimization --- speed.} Compressed modern-format images, lazy
      loading, deferred scripts, CDN delivery, targeting a load under three
      seconds on mobile data.
\item \textbf{Optimization --- checkout.} One-page guest checkout with a
      four-field form, saved payment methods and wallets, a visible progress
      indicator, and no forced account creation --- forced registration is the
      single largest abandonment cause for this persona.
\item \textbf{Optimization --- engagement and recovery.} Exit-intent offer,
      abandoned-cart e-mail and push, personalised recommendations, and
      heatmaps and session recordings to locate friction continuously.
\item \textbf{Copy suggestions.} Headline: ``Everything for your week, delivered
      by tomorrow''. Sub-line: ``Free delivery over \Rs 499. Returns in one
      tap.'' CTA: ``Shop today's deals'' rather than ``Submit''. Form microcopy:
      ``Takes 30 seconds --- no account needed''. Product copy: specific,
      benefit-led claims rather than adjectives.
\item \textbf{Measurement.} Conversion rate, mobile against desktop conversion
      gap, checkout completion rate, average order value, mobile bounce rate and
      repeat purchase rate --- with A/B tests on the hero headline and the CTA
      running continuously.
\end{apts}
\scheme{10 M. The question names four deliverables --- layout, UX and
engagement, design principles, optimisation strategies and copy --- so roughly
2--2.5 M each. Actual suggested copy strings earn marks that a description of
``good copy'' does not; write the words you would put on the page.}

\qq{c4b6}{A travel agency's landing page for a limited-time holiday offer gets
many ad clicks but converts poorly. Evaluate the possible reasons and suggest
optimizations for headlines, call-to-action buttons and visuals.}{Improvement
CIE, 28 Aug 2024 --- 10 M}
\textit{Possible reasons for the low conversion}
\begin{apts}
\item \textbf{Message mismatch} --- the advertisement promised a specific offer
      that the page does not restate, so the visitor believes they have arrived
      in the wrong place.
\item \textbf{Vague headline} that describes the agency rather than the offer.
\item \textbf{The offer and price sit below the fold,} so most visitors never
      see them.
\item \textbf{Competing or invisible calls to action} --- several links of equal
      weight, or none that is visually dominant.
\item \textbf{An over-long enquiry form} demanding details the visitor is not
      ready to give at this stage.
\item \textbf{No urgency conveyed} even though the offer is time-limited, and no
      trust signals --- no reviews, certifications or cancellation policy.
\item \textbf{Slow load and poor mobile rendering,} plus generic stock imagery
      that could belong to any agency, and hidden total cost or unclear
      inclusions. Site navigation left on the page also lets visitors wander
      away.
\end{apts}

\textit{Optimizations}
\begin{apts}\setcounter{enumi}{7}
\item \textbf{Headline.} Restate the advertisement's exact promise with
      specifics --- ``Kerala backwaters, 4 nights from \Rs 18{,}999 --- book by
      Sunday'' --- with a sub-headline listing what is included. Specificity is
      what converts on an offer page.
\item \textbf{Call-to-action buttons.} One dominant high-contrast button
      repeated above and below the fold, with active benefit copy --- ``Reserve
      my package'' rather than ``Submit'' --- plus reassurance microcopy beneath
      it such as ``Free cancellation up to 7 days''.
\item \textbf{Visuals, and everything else.} Replace stock photography with real
      destination and customer images, add a short itinerary video, and show a
      countdown timer for the deadline. Cut the form to name, phone and travel
      date; place testimonials and ratings adjacent to the CTA; strip the site
      navigation from the landing page; compress images for a sub-three-second
      mobile load; and A/B test the headline and CTA. \emph{Measure with:}
      conversion rate, form-start to completion ratio, scroll depth and session
      recordings to confirm the fix worked.
\end{apts}
\scheme{10 M. Roughly 5 M for the diagnosis and 5 M for the optimisations, and
the three elements named in the question --- headlines, CTA buttons, visuals ---
must each be addressed under their own heading, since the examiner marks against
those words.}

\qq{c4b7}{How do online applications serve as effective tools for mobile
marketing? What key features should businesses incorporate into their apps to
maximize marketing potential?}{Improvement CIE, 28 Aug 2024 --- 10 M}
\textit{Why apps are effective}
\begin{apts}
\item \textbf{An owned, unthrottled channel.} Push notifications reach the user
      directly without an algorithm deciding whether the message is delivered.
\item \textbf{Far higher engagement and conversion} than mobile web among the
      same users, because installation is itself an act of commitment.
\item \textbf{Rich first-party behavioural data} --- every screen, search and
      purchase, tied to a known user, which is increasingly valuable as
      third-party tracking degrades.
\item \textbf{Device capability access} --- GPS for location marketing, camera
      for scanning and AR try-on, biometrics for frictionless payment, wallet
      integration.
\item \textbf{Personalisation at the individual level,} driving recommendation
      engines that raise average order value.
\item \textbf{Persistent brand presence} --- the icon on the home screen is
      continuous, unpaid brand exposure.
\item \textbf{Loyalty and retention mechanics} --- points, streaks, tiers and
      wallet passes that are impractical on mobile web.
\item \textbf{Offline capability and speed,} giving a better experience on weak
      connections than any browser session.
\end{apts}

\textit{Key features to incorporate}
\begin{apts}\setcounter{enumi}{8}
\item \textbf{Acquisition and onboarding:} fast OTP or social sign-up, a short
      value-first onboarding flow, and a push-permission request made only
      \emph{after} the user has experienced value --- asking on first launch is
      how permission is permanently lost.
\item \textbf{Engagement and marketing:} segmented, deep-linked push
      notifications and in-app messaging; a personalised home feed and
      recommendation engine; in-app offers, coupons and gamification;
      referral and loyalty programmes; and social sharing.
\item \textbf{Conversion and retention:} one-tap checkout with saved payment
      and wallets; saved addresses and reorder; wishlist and price-drop alerts;
      abandoned-cart recovery; in-app support or chatbot; and ratings and
      review prompts.
\item \textbf{Measurement and hygiene:} analytics SDK with event tracking and
      attribution, A/B testing and feature flags, crash reporting, a preference
      centre for notification control, and transparent privacy and consent
      handling --- since an app that over-notifies is uninstalled, and an
      uninstalled app has no marketing value at all.
\end{apts}
\scheme{10 M. Roughly 5 M for why apps work as marketing tools and 5 M for the
feature checklist. Grouping the features by purpose --- acquisition,
engagement, conversion, measurement --- reads far better than an undifferentiated
list of eleven items.}

\qq{c4b8}{A luxury cosmetics brand is launching a flash sale exclusively through
mobile. Design a mobile marketing campaign that drives awareness of the sale and
maximizes conversion during the limited window.}{Improvement CIE, 28 Aug 2024
--- 10 M}
\textit{Pre-launch --- awareness, seven days out}
\begin{apts}
\item \textbf{Teaser push and SMS} to the segmented installed base, revealing
      the date but not the discount, building anticipation rather than
      announcing a sale.
\item \textbf{Influencer and Reels previews} with beauty creators, plus an
      in-app countdown banner on every screen.
\item \textbf{A ``notify me'' opt-in} that both captures purchase intent and, in
      the same tap, secures push permission --- the highest-value action of the
      pre-launch phase.
\item \textbf{Paid social and mobile display retargeting} to past purchasers,
      cart abandoners and lookalike audiences.
\end{apts}

\textit{Launch --- the conversion window}
\begin{apts}\setcounter{enumi}{4}
\item \textbf{Timed push at each segment's peak activity hour,} deep-linked
      straight to the sale screen rather than the home screen.
\item \textbf{App-exclusive early access} one hour before the website, which
      converts the sale into an install driver as well as a revenue event.
\item \textbf{Scarcity and urgency on-screen} --- a visible countdown timer and
      limited-stock indicators, which suit a luxury flash sale where exclusivity
      is part of the proposition.
\item \textbf{Personalised product rows} built from browsing and purchase
      history, and one-tap checkout with saved cards and wallets, since every
      additional field costs conversions during a short window.
\end{apts}

\textit{During and after}
\begin{apts}\setcounter{enumi}{8}
\item \textbf{Recovery and support:} abandoned-cart push within 20 minutes;
      geo-fenced offers with click-to-map near flagship stores; in-app live
      chat; and free-shipping threshold prompts to raise average order value.
\item \textbf{Post-sale and technical.} Thank-you and review-request push, a
      win-back offer to browsers who did not buy, and a wallet pass for the next
      drop. Technically the campaign requires sub-three-second load, server
      capacity provisioned for the traffic spike, and a fully mobile-first
      design. \emph{Measure:} push open rate, sale-screen sessions, add-to-cart
      rate, checkout completion, conversion rate, average order value, revenue
      per push, opt-out rate and installs during the window.
\end{apts}
\scheme{10 M. Structure by campaign phase --- pre-launch, launch, during, post
--- which is what earns the ``design a campaign'' marks; roughly 2.5 M per phase.
The answer must be recognisably mobile-exclusive and recognisably a
\emph{flash} sale, with urgency and scarcity mechanics explicit.}

\qq{c4b9}{You are launching a promotional campaign for a new product using mobile
marketing. Prepare a detailed plan including mobile advertising channels, push
notifications, app-based promotions and integration with social media.}{Improvement
CIE, 4 Jun 2025 --- 10 M}
\begin{apts}
\item \textbf{Objective and audience.} A SMART objective --- for example 20{,}000
      installs and 5{,}000 first purchases in eight weeks at a defined cost per
      acquisition --- and a mobile persona specifying device tier, active hours,
      preferred platforms and data-cost sensitivity.
\item \textbf{Mobile advertising --- search and intent.} Mobile search
      advertisements on high-intent queries with location and call extensions,
      and click-to-call advertisements for enquiry-driven products.
\item \textbf{Mobile advertising --- display and video.} In-app banner and
      native placements on contextually relevant apps; rewarded video for cheap
      reach; interstitials placed at natural app breakpoints rather than
      mid-task; expandable rich media for the brand story; and click-to-map
      advertisements to drive footfall.
\item \textbf{Push notifications --- the sequence.} Teaser, launch, reminder and
      last chance, each segmented by behaviour, deep-linked to the relevant
      screen, frequency-capped and A/B tested on copy and send time.
\item \textbf{Push notifications --- permission discipline.} Request permission
      only after the user has experienced value, offer a preference centre, and
      suppress the disengaged --- because an opt-out is permanent and an
      uninstall irreversible.
\item \textbf{App-based promotions.} App-exclusive early access and pricing,
      in-app coupons, gamified spin-to-win, loyalty points, referral codes,
      wallet passes and one-tap checkout.
\item \textbf{Social media integration --- organic.} Reels and Stories carrying
      the campaign hashtag, influencer seeding, user-generated content contests,
      and a messaging broadcast channel for order updates and offers.
\item \textbf{Social media integration --- paid and data.} Shoppable posts and
      product tags, click-to-WhatsApp advertisements, and retargeting audiences
      built from app and site events through the pixel and SDK --- which is what
      makes social and app one campaign rather than two.
\item \textbf{Cross-channel integration.} QR codes on packaging, print and
      in-store signage linking to the mobile landing page; one consistent
      creative and offer across every touchpoint; and mobile treated as the
      connective tissue between online and offline media.
\item \textbf{Measurement.} Installs and cost per install, push open and opt-out
      rates, in-app conversion rate, coupon redemption, ROAS, and lifetime value
      by acquisition source --- reviewed weekly, with budget reallocated to the
      channels producing profitable customers rather than cheap installs.
\end{apts}
\scheme{10 M. The question names four components explicitly --- advertising
channels, push notifications, app promotions, social integration --- so roughly
2--2.5 M each with the remainder for objectives and measurement. Answer under
those four headings.}

\qq{c4b10}{Discuss the relationship between design principles and website copy.
How can visual design and well-crafted content work together to enhance user
experience and achieve website objectives? Include examples.}{Improvement CIE,
28 Aug 2024 --- 10 M}
Design and copy are not sequential tasks --- design does not decorate finished
copy, and copy does not fill finished layouts. Each design principle has a copy
counterpart, and the page works only when both are satisfied.

\begin{apts}
\item \textbf{Visual hierarchy $\leftrightarrow$ message hierarchy.} The largest
      element must carry the most important message. A giant headline saying
      ``Welcome to our website'' is a hierarchy correctly built around the wrong
      sentence.
\item \textbf{Simplicity $\leftrightarrow$ concision.} Whitespace fails if the
      paragraph inside it runs to 200 words; clean layout and long-winded copy
      cancel each other out.
\item \textbf{Consistency $\leftrightarrow$ consistent voice and terminology.}
      If the button style is uniform but one says ``Buy now'', another ``Submit''
      and a third ``Proceed'', the interface is still inconsistent.
\item \textbf{Contrast and emphasis $\leftrightarrow$ the single call to
      action.} The one visually dominant element and the one action the copy
      asks for must be the same element.
\item \textbf{Readability $\leftrightarrow$ plain language.} Type size, line
      length and contrast govern whether text can be read; sentence length and
      jargon govern whether it can be understood. Failing either loses the
      reader identically.
\item \textbf{Imagery $\leftrightarrow$ specific claims.} A photograph shows
      what the product looks like; the copy supplies the number, the guarantee
      and the proof. Neither persuades alone.
\item \textbf{Whitespace and rhythm $\leftrightarrow$ scannable structure.}
      Subheadings, bullets and short paragraphs are simultaneously a design
      decision and a writing decision.
\end{apts}

\textit{Worked examples of successful integration}
\begin{apts}\setcounter{enumi}{7}
\item \textbf{Hero section.} A high-contrast headline stating one benefit, a
      supporting line, one accent-coloured button, over an image showing the
      product in use --- the eye reaches the headline first, the copy answers the
      question it raises, and the button is where the eye lands next.
\item \textbf{Pricing table.} Three aligned columns with the recommended tier
      visually elevated, copy naming who each tier is for --- ``for teams of
      2--10'' --- rather than merely listing features. Design directs the choice;
      copy justifies it.
\item \textbf{Failure modes, for contrast.} A beautiful layout carrying vague
      copy --- ``innovative solutions for a connected world'' --- communicates
      nothing; and strong specific copy set in low-contrast 12-pixel grey type
      is never read at all. Both fail, and both fail for the same reason: one
      half of the pairing was treated as decoration.
\end{apts}
\scheme{10 M. Roughly 6 M for the principle-by-principle relationship --- a
two-column table is the most efficient presentation --- and 4 M for the worked
examples, which the question demands. Including a failure example as well as
successes usually earns the top band.}

\qq{c4b11}{What is the role of user interface and user experience design in web
design for marketing? How can businesses ensure a positive user journey on their
websites?}{SEE, Oct/Nov 2023, Q7a --- 8 M}
\textbf{UI} is what the user sees and interacts with --- layout, colour,
typography, buttons, icons, spacing, visual states. \textbf{UX} is the whole
experience of using the site --- research, information architecture, flows,
usability, accessibility and satisfaction. UI is a component of UX: a beautiful
interface over a broken flow is good UI and bad UX.

\textit{Their role in marketing}
\begin{apts}
\item \textbf{They convert the traffic marketing pays for.} Acquisition spend is
      realised only at the destination, so UI/UX determines the return on every
      other marketing rupee.
\item \textbf{They establish credibility within seconds,} which for an unknown
      brand decides whether anything else is read.
\item \textbf{They reduce friction and therefore abandonment,} particularly in
      forms and checkout.
\item \textbf{They support SEO indirectly,} because engagement, speed and mobile
      usability are ranking signals.
\item \textbf{They differentiate the brand} where products are otherwise
      comparable, and drive retention and repeat visits.
\end{apts}

\textit{Ensuring a positive user journey}
\begin{apts}\setcounter{enumi}{5}
\item \textbf{Research the users and map the journey} --- personas, tasks and
      touchpoints --- before designing anything.
\item \textbf{Design mobile-first and responsively,} and keep load times under
      three seconds.
\item \textbf{Make navigation and information architecture obvious,} with
      search, breadcrumbs and a shallow hierarchy.
\item \textbf{Maintain message match and consistency} from advertisement to
      landing page to checkout, so the journey never contradicts itself.
\item \textbf{Minimise friction} --- fewest form fields, guest checkout,
      progress indicators, clear error messages that say how to fix the problem.
\item \textbf{Meet accessibility standards} for contrast, keyboard operation and
      screen readers, so the journey exists for everyone.
\item \textbf{Test with real users and iterate} --- usability testing, heatmaps,
      session recordings and A/B tests --- and measure the journey with bounce,
      completion, task-success and conversion rates.
\end{apts}
\scheme{8 M. Roughly 2 M for defining UI and UX and their relationship, 3 M for
their role in marketing, and 3 M for ensuring a positive journey. Confusing UI
with UX, or treating them as synonyms, costs the first two marks immediately.}

\qq{c4b12}{Describe the importance of testing and optimization in web design for
marketing. What tools and methods can be used to analyze and improve website
performance?}{SEE, Oct/Nov 2023, Q7b --- 8 M}
\textit{Importance}
\begin{apts}
\item \textbf{Design opinion is unreliable.} What a team believes will convert
      and what does convert diverge often enough that testing, not judgement, is
      the only dependable method.
\item \textbf{It raises revenue without raising traffic cost,} improving the
      economics of every acquisition channel at once.
\item \textbf{It reduces the risk of redesigns,} since changes are validated on
      a fraction of traffic before full rollout.
\item \textbf{It catches defects that analytics alone hides} --- a broken form
      on one browser, a checkout failure on one device --- which can silently
      cost a large share of revenue.
\item \textbf{It creates a culture of continuous improvement,} where the site is
      treated as a product under development rather than a project that ended.
\end{apts}

\textit{Methods and their tools}
\begin{apts}\setcounter{enumi}{5}
\item \textbf{Quantitative analysis} --- Google Analytics 4 for traffic,
      funnels, segments and conversions; Google Search Console for search
      performance.
\item \textbf{Behavioural analysis} --- Hotjar or Microsoft Clarity for
      heatmaps, scroll maps, session recordings and form analytics; on-site
      surveys for voice of customer.
\item \textbf{Performance testing} --- PageSpeed Insights, Lighthouse, GTmetrix
      and WebPageTest for load time and Core Web Vitals.
\item \textbf{Experimentation} --- A/B and multivariate testing through
      Optimizely, VWO or AB Tasty, run to statistical significance.
\item \textbf{Usability and accessibility testing} --- moderated and unmoderated
      user testing, five-second tests, WAVE and axe for accessibility; plus
      cross-browser and cross-device testing (BrowserStack) and technical
      auditing (Screaming Frog, SEMrush Site Audit).
\end{apts}
\scheme{8 M. Roughly 4 M for the importance and 4 M for methods paired with
named tools. Naming actual tools rather than saying ``analytics software''
carries marks; a table of method against tool against what it reveals is the
most efficient format.}

\qq{c4b13}{Discuss push notifications in mobile marketing. How can businesses use
them to engage and retain mobile app users? Provide examples of effective push
notification strategies.}{SEE, Oct/Nov 2023, Q8a --- 8 M}
\textbf{Definition.} Push notifications are short messages sent by an
application directly to a user's device screen, appearing whether or not the app
is open. They require explicit permission, and are the only marketing channel
that reaches an installed user without an intermediary algorithm deciding
delivery.

\textit{Uses for engagement and retention}
\begin{apts}
\item \textbf{Re-engagement} of users who have not opened the app recently,
      before dormancy becomes uninstallation.
\item \textbf{Transactional and status updates} --- order confirmed, shipped,
      out for delivery --- which are opened almost universally and build the
      habit of trusting notifications from that app.
\item \textbf{Cart and browse abandonment recovery,} typically the highest-value
      automated push available.
\item \textbf{Personalised offers and recommendations} based on past behaviour.
\item \textbf{Time-sensitive promotions} --- flash sales, price drops,
      back-in-stock and low-stock alerts.
\item \textbf{Location-triggered messages} when a user enters a geofence near a
      store.
\item \textbf{Content, streaks and loyalty prompts} that build routine usage.
\end{apts}

\textit{Effective strategies, with examples}
\begin{apts}\setcounter{enumi}{7}
\item \textbf{Ask permission at the right moment} --- after the user has
      experienced value, with an explanation of what they will receive, never on
      first launch.
\item \textbf{Segment and personalise} --- send to behaviour-defined groups, not
      the whole base. \emph{Example:} a cart reminder 20 minutes after
      abandonment, naming the item left behind.
\item \textbf{Time by behaviour, not by convenience} --- send at each segment's
      observed active hour and respect time zones. \emph{Example:} a food
      delivery app pushing at 12:30 pm and 8 pm rather than mid-morning.
\item \textbf{Be specific, brief and deep-linked} --- state the benefit in under
      ten words and land the tap on the exact screen. \emph{Example:} ``Your
      size is back in stock --- 3 left'' opening the product page.
\item \textbf{Cap frequency, give a preference centre, and A/B test everything}
      --- copy, emoji, timing and offer --- measuring opt-out rate alongside open
      rate. \emph{Example:} a 30-day win-back offering a specific incentive
      rather than a generic ``we miss you''. Over-notification is the single
      largest cause of uninstalls, so restraint is a strategy, not a compromise.
\end{apts}
\scheme{8 M. Roughly 1 M for the definition, 3--4 M for the engagement and
retention uses, and 3--4 M for strategies with concrete examples --- the
question asks for examples explicitly, so generic advice without them loses
those marks.}

\qq{c4b14}{Explain the significance of mobile analytics in mobile marketing
campaigns. What key metrics should marketers track to measure success?}{SEE,
Oct/Nov 2023, Q8b --- 8 M}
\textit{Significance}
\begin{apts}
\item \textbf{It measures a journey that web analytics cannot see.} App
      sessions, screens and in-app events are invisible to conventional web
      tracking, so without mobile analytics a large part of the customer
      journey is simply unrecorded.
\item \textbf{It attributes installs and revenue to source,} which is the only
      way to know which campaign actually produced paying users rather than
      cheap installs.
\item \textbf{It exposes the retention problem.} Most installed apps are
      abandoned quickly; only cohort retention analysis reveals whether
      acquisition spend is building a user base or filling a leaking bucket.
\item \textbf{It supports personalisation and segmentation} from behavioural
      data, and enables in-app A/B testing and feature experimentation.
\item \textbf{It surfaces technical failures} --- crashes, slow screens, failed
      payments --- that destroy conversion silently and would otherwise be
      reported only by the few users who complain.
\end{apts}

\textit{Key metrics}
\begin{apts}\setcounter{enumi}{5}
\item \textbf{Acquisition:} installs, cost per install, install source and
      campaign attribution, click-to-install rate.
\item \textbf{Activation and engagement:} onboarding completion rate, daily and
      monthly active users, DAU/MAU stickiness ratio, sessions per user, session
      length, screens per session, feature adoption.
\item \textbf{Retention:} Day 1, Day 7 and Day 30 retention, cohort retention
      curves, churn rate and uninstall rate.
\item \textbf{Conversion and revenue:} in-app conversion rate, funnel completion,
      average revenue per user, average order value, lifetime value, and the
      LTV-to-CAC ratio --- which is the metric that decides whether the whole
      campaign is viable.
\item \textbf{Channel and technical health:} push opt-in rate, push open rate
      and opt-out rate, deep-link performance, crash-free session rate, screen
      load time, and app-store rating and review sentiment.
\end{apts}
\scheme{8 M. Roughly 4 M for the significance and 4 M for the metrics, best
grouped by funnel stage --- acquisition, engagement, retention, revenue ---
rather than listed flat. Naming LTV against CAC, and retention cohorts, marks
the answer as informed rather than generic.}

\end{qlist}
